%! AP Statistics — Unit 2, Lesson 2.8 — Homework — Answer Key
\documentclass[10pt]{article}
\usepackage{apstats-article}
\usepackage{apstats-key}

\begin{document}

\pageheader{Unit 2, Lesson 2.8}{Homework --- Answer Key}
\namedateperiod

\vspace{0.14in}

% ── PROBLEM 1 ─────────────────────────────────────────────────────────────────
\begin{scenariobox}[Context A: Study Hours and GPA]{sky}
\small A researcher records weekly study hours ($x$) and semester GPA ($y$) for 50 college
students. Summary statistics: $\bar{x} = 18$ hr, $\bar{y} = 3.1$, $s_x = 6$ hr, $s_y = 0.5$,
$r = 0.70$. The data ranged from 5 to 35 hours per week.
\end{scenariobox}

\vspace{0.1in}

\begin{enumerate}[leftmargin=1.6em, itemsep=0.18in]

  \item Compute $b$ and $a$. Show work.\\[8pt]
        \ansline{$b = r\,\frac{s_y}{s_x} = 0.70\cdot\frac{0.5}{6} \approx 0.0583$ GPA points per hr}\\[3pt]
        \ansline{$a = \bar{y} - b\bar{x} = 3.1 - (0.0583)(18) = 3.1 - 1.05 \approx 2.05$}

  \item Write the regression equation using variable names.\\[8pt]
        \ansline{$\widehat{\text{GPA}} = 2.05 + 0.058(\text{study hours})$}

  \item Interpret the \textbf{slope} in context.\\[8pt]
        \ansline{For each additional hour of weekly study time, the predicted semester GPA increases}\\[3pt]
        \ansline{by approximately 0.058 points.}

  \item Interpret the \textbf{$y$-intercept} in context. Is it meaningful? Explain.\\[8pt]
        \ansline{When a student studies 0 hours per week, the predicted GPA is approximately 2.05.}\\[3pt]
        \ansline{This \textbf{is meaningful in context} because 0 hours per week is within the range (5--35 hr is}\\[3pt]
        \ansline{the data range, so 0 is technically outside; state it is at the boundary / barely outside).}

  \item Compute $r^2$. Write the full AP-template interpretation.\\[8pt]
        \ansline{$r^2 = (0.70)^2 = 0.49$}\\[3pt]
        \ansline{About 49\% of the variation in semester GPA is explained by the least-squares}\\[3pt]
        \ansline{regression of GPA on weekly study hours.}

  \item Predict the GPA for a student who studies 25 hours per week. Classify as interpolation
        or extrapolation.\\[8pt]
        \ansline{$\widehat{\text{GPA}} = 2.05 + 0.058(25) = 2.05 + 1.45 = 3.50$}\\[3pt]
        \ansline{\textbf{Interpolation} --- 25 hours is within the data range (5--35 hr).}

\end{enumerate}

\vspace{0.18in}

% ── PROBLEM 2 ─────────────────────────────────────────────────────────────────
\begin{scenariobox}[Context B: Altitude and Temperature]{goldbg}
\small A meteorologist measures altitude in thousands of feet ($x$) and air temperature in
\textdegree F ($y$) at 40 locations. Summary statistics: $\bar{x} = 5$ (thousand ft),
$\bar{y} = 42$\textdegree F, $s_x = 2$ (thousand ft), $s_y = 14$\textdegree F, $r = -0.90$.
\end{scenariobox}

\vspace{0.1in}

\begin{enumerate}[leftmargin=1.6em, itemsep=0.18in, start=7]

  \item Compute $b$, $a$, and $r^2$. Write the regression equation.\\[8pt]
        \ansline{$b = (-0.90)\cdot\frac{14}{2} = -6.3$\textdegree F per thousand ft; \quad $a = 42 - (-6.3)(5) = 42 + 31.5 = 73.5$}\\[3pt]
        \ansline{$r^2 = (-0.90)^2 = 0.81$; \quad $\widehat{\text{temp}} = 73.5 - 6.3(\text{altitude})$}

  \item Interpret $r^2$ in context.\\[8pt]
        \ansline{About 81\% of the variation in air temperature is explained by the least-squares}\\[3pt]
        \ansline{regression of air temperature on altitude.}

  \item A data point at altitude 18,000 ft ($x = 18$) with a temperature of $-65$\textdegree F
        is discovered. Is this likely to be an influential point? Why?\\[8pt]
        \ansline{Yes. $x = 18$ is far outside the data range, giving this point \textbf{high leverage}.}\\[3pt]
        \ansline{However, the point follows the negative trend (higher altitude $\to$ lower temperature),}\\[3pt]
        \ansline{so it may increase $r^2$ while pulling the slope downward. It is influential.}

\end{enumerate}

\vspace{0.18in}

% ── EXTENSION ────────────────────────────────────────────────────────────────
\begin{extensionbox}
\small\textbf{Extension --- optional}

\vspace{4pt}
For Context B, the point $(18, -65)$ is added. The new regression equation becomes
$\widehat{\text{temp}} = 82 - 8.5(\text{altitude})$ and $r^2 = 0.94$.

\begin{enumerate}[leftmargin=1.6em, itemsep=6pt]
  \item \ans{Slope: $-6.3 \to -8.5$ (more negative --- the outlier pulls the line steeper). Intercept: $73.5 \to 82$ (increased to match the steeper slope through $(\bar{x},\bar{y})$). $r^2$: $0.81 \to 0.94$ (increased --- the outlier follows the trend and extends the range, strengthening the apparent linear fit).}
  \item \ans{Yes, this confirms influence: the slope, intercept, and $r^2$ all changed substantially when the point was added.}
  \item \ans{Do not remove it automatically. First verify the measurement (could 18,000 ft data be real?). If so, report it as an influential point that follows the trend and note its effect on the model.}
\end{enumerate}
\end{extensionbox}

\vspace{0.16in}

% ── PREVIEW ──────────────────────────────────────────────────────────────────
\begin{tcolorbox}[colback=greenbg, colframe=greenacc, boxrule=0.8pt, arc=2mm,
    left=4mm, right=4mm, top=2.5mm, bottom=2.5mm]
\small\textbf{Preview for Lesson 2.9:} The LSRL is the best \emph{linear} model --- but what
if the relationship is curved? When a residual plot shows a pattern (not random scatter), a
\textbf{transformation} (log, square root, power) can straighten the data before applying
regression. We'll explore when and how to transform variables next class.
\end{tcolorbox}

\begin{teachernote}
Problem 1 item 4 ($y$-intercept meaningfulness): $x = 0$ hr is just outside the data range
(5--35 hr). Accept either ``meaningful'' or ``not quite meaningful'' if well-justified;
the key point is that students engage with the question rather than skip it. Award full credit
for acknowledging the boundary issue.

Problem 2 extension: the point follows the trend but has extreme leverage, so $r^2$ increases.
This is the counterintuitive case --- an influential point that \emph{improves} the apparent fit.
Use it to reinforce: high leverage + trend-consistent $\Rightarrow$ possibly inflated $r^2$.
\end{teachernote}

\end{document}
